%! Sinusoidal Function Context and Data Modeling — Unit 3, Lesson 3.7 — Guided Notes (Answer Key)
\documentclass[10pt]{article}
\usepackage{precalculus-article}
\usepackage{precalculus-key}   % pulls in -boxes; do NOT also load -boxes
\newcommand{\TallMath}[1]{$\displaystyle #1\rule[-1.4em]{0pt}{3.2em}$}

\begin{document}

\pageheader{Unit 3, Lesson 3.7}{Guided Notes --- Answer Key}
\namedateperiod

\vspace{0.10in}

%----------------------------------------------------------------------
% Objectives
%----------------------------------------------------------------------
\begin{objectivebox}
\textbf{I will be able to\ldots}
\begin{enumerate}[leftmargin=1.5em, itemsep=4pt]
  \item \textbf{LO 3.7.A} --- Construct a sinusoidal model from a data set by
    computing amplitude, midline, period, and phase shift, and validate the
    model against known data values.
    \ansline{Construct sinusoidal models using amplitude, midline, period, and phase shift.}
  \item \textbf{LO 3.7.A} --- Use a sinusoidal model to predict output values
    and identify the contextual domain over which the model is meaningful.
    \ansline{Use sinusoidal models to predict values over the contextual domain.}
\end{enumerate}
\end{objectivebox}

\vspace{0.08in}

%----------------------------------------------------------------------
% Vocabulary
%----------------------------------------------------------------------
\begin{vocabbox}
\begin{description}[leftmargin=0pt, itemsep=6pt]
  \item[\termblanklong{sinusoidal model}]
    A function of the form $f(x) = A\sin(B(x - C)) + D$ (or cosine form) used
    to represent a \ans{periodic} real-world quantity.

  \item[\termblanklong{amplitude}]
    $A = \dfrac{\max - \min}{2}$; measures how far the function deviates from
    its \ans{midline}.

  \item[\termblanklong{vertical shift / midline}]
    $D = \dfrac{\max + \min}{2}$; the horizontal line $y = D$ about which the
    function \ans{oscillates}.

  \item[\termblanklong{period}]
    The \ans{smallest} input-value interval between consecutive maxima
    (or minima) in the data. Equals $\dfrac{2\pi}{B}$.

  \item[\termblanklong{phase shift}]
    $C$; shifts the graph \ans{right} when $C > 0$. Determined
    by comparing a data maximum to the standard model maximum location.

  \item[\termblanklong{contextual domain}]
    The set of input values over which the model is \ans{meaningful / valid}
    given the real-world context.
\end{description}
\end{vocabbox}

\vspace{0.08in}

%----------------------------------------------------------------------
% Hook
%----------------------------------------------------------------------
\begin{hookbox}
\textbf{Entry Question:} The table below shows average monthly high
temperatures (${}^\circ$F) for a city over one year.

\vspace{4pt}
\begin{center}
\renewcommand{\arraystretch}{1.3}
\small
\begin{tabular}{|c|c|c|c|c|c|c|c|c|c|c|c|c|}
\hline
\textbf{Month} $x$ & 1 & 2 & 3 & 4 & 5 & 6 & 7 & 8 & 9 & 10 & 11 & 12 \\
\hline
\textbf{Temp} $T(x)$ (${}^\circ$F) & 38 & 42 & 54 & 65 & 74 & 82 & 86 & 84 & 75 & 63 & 49 & 40 \\
\hline
\end{tabular}
\end{center}

\vspace{6pt}
\textbf{Q1.} Which month has the maximum temperature? \ans{7 (July)}
\quad Which has the minimum? \ans{1 (January)}

\vspace{4pt}
\textbf{Q2.} What kind of function could model this data?
\ansline{A sinusoidal (sine or cosine) function, because the data rises, falls, and would repeat periodically over successive years.}
\end{hookbox}

\vspace{0.08in}

%----------------------------------------------------------------------
% LO 3.7.A — Four-Step Model Building Process
%----------------------------------------------------------------------
\begin{notesbox}{LO 3.7.A --- Building a Sinusoidal Model: Four Steps}

We use the model $T(x) = A\sin(B(x - C)) + D$ where:
\begin{itemize}[leftmargin=1.4em, itemsep=2pt]
  \item $A =$ amplitude, \quad $D =$ midline (vertical shift),
  \item $B = \dfrac{2\pi}{k}$ ($k =$ period), \quad $C =$ phase shift.
\end{itemize}

\vspace{8pt}
\textbf{Step 1 --- Find the Period $k$}

The period equals the input-value distance between consecutive \ans{maxima}
(or consecutive \ans{minima}).

In the temperature table, consecutive maxima occur at months \ans{7}
and \ans{7} (next year), so $k \approx$ \ans{12}.

$B = \dfrac{2\pi}{k} = $ \ans{$\dfrac{2\pi}{12} = \dfrac{\pi}{6} \approx 0.524$}

\vspace{10pt}
\textbf{Step 2 --- Compute Amplitude and Midline}

$\max = $ \ans{86} \quad $\min = $ \ans{38}

$A = \dfrac{\max - \min}{2} = \dfrac{{\color{keyred}\mathbf{86}} - {\color{keyred}\mathbf{38}}}{2} =$ \ans{24}

$D = \dfrac{\max + \min}{2} = \dfrac{{\color{keyred}\mathbf{86}} + {\color{keyred}\mathbf{38}}}{2} =$ \ans{62}

The midline represents \ans{the average of the hottest and coolest monthly temperatures ($62^\circ$F)} in context.

\vspace{10pt}
\textbf{Step 3 --- Determine the Phase Shift $C$}

For $T(x) = A\sin(B(x - C)) + D$, the maximum occurs when $B(x - C) = \dfrac{\pi}{2}$,
so $x = C + \dfrac{\pi}{2B}$.

The data maximum occurs at $x = $ \ans{7}. Set equal and solve for $C$:

\vspace{4pt}
$C + \dfrac{\pi}{2B} = $ \ans{7} \quad $\Rightarrow$ \quad $C \approx$ \ans{$7 - \dfrac{\pi}{2 \cdot \pi/6} = 7 - 3 = 4$}

\vspace{6pt}
\textit{Tip:} If you use a \textbf{cosine} model $T(x) = A\cos(B(x-C)) + D$, the maximum
occurs when $B(x - C) = 0$, so $C = x_{\max}$ \ans{$= 7$ directly}.

\vspace{10pt}
\textbf{Step 4 --- Write and Validate the Model}

$T(x) = $ \ans{$24\sin\!\left(\tfrac{\pi}{6}(x - 4)\right) + 62$}

\vspace{6pt}
Validate: substitute a known data point (e.g., $x = 1$, $T = 38$):

$T(1) \approx$ \ans{$24\sin(-\tfrac{\pi}{2}) + 62 = 24(-1)+62 = 38$} \quad Residual $= $ \ans{$38 - 38 = 0$}

\vspace{6pt}
\textbf{Contextual Domain:} This model is valid for \ans{$1 \leq x \leq 12$ (one calendar year for this location)}.

\end{notesbox}

\vspace{0.08in}

%----------------------------------------------------------------------
% Practice
%----------------------------------------------------------------------
\begin{practicebox}

\textbf{Guided Practice.}\quad The table below shows the depth of water (in feet)
at a harbor entrance over one tidal cycle (12 hours).

\vspace{4pt}
\begin{center}
\renewcommand{\arraystretch}{1.35}
\begin{tabular}{|c|c|c|c|c|c|c|c|}
\hline
Time $t$ (hr) & 0 & 2 & 4 & 6 & 8 & 10 & 12 \\
\hline
Depth $d(t)$ (ft) & 5 & 9.5 & 13 & 9.5 & 5 & 0.5 & 5 \\
\hline
\end{tabular}
\end{center}

\vspace{4pt}
(a)\quad Identify $\max = $ \ans{13} at $t = $ \ans{4}
and $\min = $ \ans{0.5} at $t = $ \ans{10}.

\vspace{4pt}
(b)\quad Compute: $A = $ \ans{$\dfrac{13-0.5}{2} = 6.25$} \quad $D = $ \ans{$\dfrac{13+0.5}{2} = 6.75$}

\vspace{4pt}
(c)\quad Estimate the period $k = $ \ans{12} hr. \quad Then $B = \dfrac{2\pi}{k} = $ \ans{$\dfrac{\pi}{6} \approx 0.524$}.

\vspace{4pt}
(d)\quad Identify the phase shift $C$, using a cosine model with maximum at $t = 4$. $C = $ \ans{4}.

\vspace{4pt}
(e)\quad Write the model: $d(t) = $ \ans{$6.25\cos\!\left(\tfrac{\pi}{6}(t - 4)\right) + 6.75$}

\vspace{4pt}
(f)\quad Use your model to predict the depth at $t = 3$.
\ansline{$d(3) = 6.25\cos(\tfrac{\pi}{6}(3-4)) + 6.75 = 6.25\cos(-\tfrac{\pi}{6})+6.75 \approx 6.25(0.866)+6.75 \approx 12.16$ ft}

\vspace{4pt}
(g)\quad Is $t = 15$ hr inside the contextual domain for one tidal cycle? Explain.
\ansline{No. The contextual domain is $0 \leq t \leq 12$ (one tidal cycle). $t = 15$ is outside that range, though the model could extend if the tidal pattern continues.}

\end{practicebox}

\end{document}
